% Summarize the research papers 12 to 15 journal papers from recent 3 years.
% Ensure the impact factor of the journal is more than 3. Write the gaps in the
% literature.

% =========================================================================== %
Modern day misinformation need not be only altered images. It can also be the
captions attached to unaltered image. \citet{sagar2026factfake} aims to verify
the efficacy of the use of IMD techniques in the context of “Multimodal
Misinformation”. This is achieved in a two-stage process starting with a MCTS
multimodal LLM supported by web search and similar RAG tools, and multiple
rounds of “Multi-Agent Debate” between a sceptic agent and a supporter agent.
This new architecture is tested against five IMD techniques, namely, NPR
\cite{tan2023npr}, BNext-M \cite{romeo2024bnextm}, GRAMNet
\cite{liu2020gramnet}, ProDet \cite{cheng2024prodet}, and
FatFormer \cite{liu2023fatformer}.
These are then benchmarked on the MMFakeBench and the DGM4 datasets.

In \citet{sun2025forgerysleuth}, the authors address limitations of using
multimodal large language models (M-LLMs) to perform image
manipulation detection (IMD) tasks by proposing “ForgerySleuth”, a
tool that aims to provide a much-detailed explanation of detected
clues and segments the regions that have been tampered. They have
also constructed a more diverse dataset, aptly titled
“ForgeryAnalysis Dataset” from existing datasets including CASIA
\cite{dong2013casia}, AutoSplice \cite{jia2023autosplice},
DEFACTO \cite{mahfoudi2019defacto}, and MIML \cite{qu2024miml}.

While the field of image synthesis with the use of artificial
intelligence has been rapidly evolving thanks to the diffusion-based
generative models, identification of such content becomes crucial in
image forensics. To address the increased existence of GAN produced
images, \citet{tran2025diffcor} propose DiffCoR. This detection
method uses Stable Diffusion Processing (SDP) module, which uses a
pre-trained model to reconstruct images using reverse diffusion and
Image Representation Learning (IRL) that applies Self-Supervised
Learning (SSL) with Latency Consistency Loss (LCL). Additionally,
Discrete Fourier Transform (DFT) allows a frequency domain analysis
for enhancing the detection of any manipulation. The dataset
introduced for this research is “ForensicsImage”, which constitutes
of both real and AI generated images from LSUN-Bedroom, CelebA-HQ and
CelebDFv2, among other diffusion models like LDM
\cite{rombach2022ldm}, DALL-E  \cite{ramesh2022dalle}, Imagen
\cite{saharia2022imagen}, and Midjourney.

In contrast to the deep learning framework, DiffCoR suggests,
\citet{bammey2024synthbuster} brings forward the idea of a
computationally lean, interpretable detector optimized for
photorealistic diffusion outputs from models like Stable Diffusion
1.3-XL \cite{rombach2022ldm}, Midjourney v5, DALL·E 2/3
\cite{ramesh2022dalle}, Glide, and Adobe Firefly, using a
cross-difference high-pass filter to highlight frequency artifacts at
periods 2, 4, and 8 in the Fourier transform of a residual image,
allowing a clear distinction of a real image from a fake one. In
addition, it extracts 135 peak magnitudes per image as features for a
lightweight HBGB classifier and handles JPEG compression through
per-quality-factor models. The research used a self-constructed
dataset from which the synthetic images were generated from text
prompts, while the real images were borrowed from the Raise-1k
dataset, a subset of the Raise \cite{dang-nguyen2015raise}
dataset. The architecture brings a good generalization to the table
but focuses only on diffusion models.

While methods such as those proposed by \citet{bammey2024synthbuster}
work effectively on high-frequency artifacts, these details are at
risk of being lost during post-processing like image compression. For
detection to work even after such alterations and to expand detection
beyond just diffusion models, \citet{manisha2025camera} propose an
approach based on source camera fingerprinting. Fingerprints
(artifacts and other information) left on photographs by physical
cameras and AI models are very different; by focusing on
mid-frequency details, this difference can be discerned, even after
compression. Since each AI model will have its unique fingerprint,
this methodology also makes it possible to identify the exact model
that has generated an image. A ResNet101 model is trained on QUFVD
images resized to 224x224 pixels. The resizing strips high frequency
data and forces the model to learn mid-frequency fingerprints. The
trained ResNet101 outputs a 2048-dimensional fingerprint, which is
then fed into the SVM, for a set of real (FFHQ) and fake (DH3)
images, to build the classification model. On the DF3 and CNN-aug
datasets, the model has an average AUC of 0.994 and showcases high
resilience to image alterations. However, the reliance on
fingerprints compromises the accuracy of the model when there are
localized image changes such as face blending.

Building upon the need to capture localized artifacts that methods
like \citet{manisha2025camera} struggle with while still making use of
frequency-domain analysis, \citet{yan2026dff} propose a Dual Frequency Branch
Framework combined with Reconstructed Sliding Windows Attention
(RSWAttention). Unlike earlier single-frequency approaches that could
lose data during compression or resizing, this architecture extracts
many perspective features by fusing the Discrete Wavelet Transform
(DWT) for local high-frequency and low-frequency textures, with the
phase component of the Fast Fourier Transform (FFT). To isolate and
find subtle tampering, RSWAttention restricts feature extraction to
the local window, which forces the model to evaluate how the
neighboring pixels are interdependent. It was evaluated against 65
distinct GAN and diffusion models, and showed good generalization,
better than baselines such as FatFormer by \citet{liu2023fatformer}. Despite the
local attention mechanism, the framework has vulnerabilities with source
fingerprinting while detecting highly localized face-swapping alterations.
